\chapter{Statistical Reporting Examples}
\label{chap:presentation}
This chapter illustrates statistical presentation with tables, figures,
algorithms, code, output, and reference formats. Replace the demonstration
content with your research while retaining the patterns you need.


\section{Writing style}
Follow \citet{schifano2026statistical} when preparing the thesis. Their guide,
\emph{Statistical Writing}, is available at
\url{https://statds.github.io/stat-writing/}. Its chapters on writing tools and
manuscript sections provide guidance on source preparation, mathematical
exposition, and statistical reporting. University requirements govern the
document layout and front matter.


Introduce the research problem, explain the gap in existing work, and state the
contribution. Define notation consistently and state assumptions before
presenting theoretical results. The abstract should stand on its own, summarize
the contribution, and avoid citations. Discuss the meaning of results rather
than repeating every value from tables and figures.


Treat equations as parts of sentences and punctuate them accordingly. Before
a displayed equation, ``is'' is sufficient; ``is given by'' is unnecessary. Use
descriptive labels and refer to equations with \verb|\eqref|. Refer to each
table and figure in the text; give each an informative, self-contained caption.
Use present tense for exposition and past tense for completed experiments.
Treat ``data'' as plural and reserve ``significant'' for its statistical
meaning.


\section{Mathematical statements and proofs}
\begin{definition}
  For $n>0$ observations $x_1,\ldots,x_n\in\RR$, the sample mean is
  \begin{equation}
    \label{eq:sample-mean}
    \bar{x} = \frac{1}{n} \sum_{i=1}^{n} x_i.
  \end{equation}
\end{definition}
\begin{theorem}
  \label{thm:constant-mean}
  If every observation equals $c$, then $\bar{x}=c$.
\end{theorem}
\begin{proof}
  Substitution in Equation~\eqref{eq:sample-mean} gives $nc/n=c$.
\end{proof}
\begin{corollary}
  The sample mean of a sequence of zeros is zero.
\end{corollary}
\begin{proof}
  Apply Theorem~\ref{thm:constant-mean} with $c=0$.
\end{proof}
Equation~\eqref{eq:centered-zero} shows that the centered observations sum to
zero. An aligned derivation gives:
\begin{align}
  \sum_{i=1}^{n}(x_i-\bar{x}) &= \sum_{i=1}^{n} x_i - n \bar{x}
  \notag\\
  &= 0.\label{eq:centered-zero}
\end{align}


\section{Differentiation and integration}
A probability density with two cases is
\begin{equation}
  \label{eq:piecewise-density}
  f(x) =
  \begin{cases}
    2x, & 0<x<1, \\
    0, & \text{otherwise}.
  \end{cases}
\end{equation}
The comma after the first condition separates the cases; the period after
``otherwise'' ends the sentence. Keep this punctuation inside the cases
rather than adding another period after the closing brace.


Use \verb|\dd| from the \texttt{physics} package for the upright differential
operator in derivatives and integration measures. For a random variable
with density $f(x)=2x$ on $0<x<1$ and zero elsewhere, normalization gives
\[
  \int_0^1 f(x)\dd{x}=\int_0^1 2x\dd{x}=1.
\]
Its expectation is obtained by integrating $x$ against the density:
\[
  \mathbb{E}(X)=\int_0^1 x f(x)\dd{x}=\frac{2}{3}.
\]
Within the support, the logarithmic derivative is
\[
  \frac{\dd}{\dd x}\log f(x)=\frac{1}{x}.
\]


\section{Lists}
\begin{itemize}
  \item Describe the research question.
  \item Explain the method and assumptions.
  \item Present results and their limitations.
\end{itemize}
\begin{enumerate}
  \item Read the observations.
  \item Calculate their mean.
  \item Check the calculation against a known example.
\end{enumerate}
\begin{description}
  \item[Input] A nonempty sequence of observations.
  \item[Output] Its arithmetic mean.
\end{description}


\section{Tables and figures}
For figure and table style tips, follow the guidance in
\emph{Statistical Writing} \citep{schifano2026statistical}, available at
\url{https://statds.github.io/stat-writing/}. Use that guide when preparing
plots, formatting tables, and writing captions for your thesis.


Table captions conventionally precede the table; figure captions follow the
illustration. Put labels after captions and refer to them by label.
Table~\ref{tab:sample-data} summarizes the values plotted in
Figure~\ref{fig:sample-data}.
\begin{table}[tbp]
  \centering
  \caption{Three observations, $x_1 = 1$, $x_2 = 2$, and $x_3 = 3$,
  and their arithmetic mean.}\label{tab:sample-data}
  \begin{tabular}{lr}
    \toprule
    Quantity & Value \\
    \midrule
    First observation & 1 \\
    Second observation & 2 \\
    Third observation & 3 \\
    Sample mean & 2 \\
    \bottomrule
  \end{tabular}
\end{table}
\begin{figure}[tbp]
  \centering
  \begin{picture}(80,40)
    \put(10,8){\vector(1,0){60}}
    \put(10,8){\vector(0,1){28}}
    \put(70,3){\makebox(0,0){$i$}}
    \put(5,35){\makebox(0,0){$x_i$}}
    \put(25,16){\circle*{2}}
    \put(40,24){\circle*{2}}
    \put(55,32){\circle*{2}}
    \put(9,16){\line(1,0){2}}
    \put(9,24){\line(1,0){2}}
    \put(9,32){\line(1,0){2}}
    \put(6,16){\makebox(0,0){1}}
    \put(6,24){\makebox(0,0){2}}
    \put(6,32){\makebox(0,0){3}}
    \put(25,4){\makebox(0,0){1}}
    \put(40,4){\makebox(0,0){2}}
    \put(55,4){\makebox(0,0){3}}
  \end{picture}
  \caption{Observation values $x_i = i$ for indices $i = 1, 2, 3$.
  The arithmetic mean of these values is 2.}
  \label{fig:sample-data}
\end{figure}
Figure~\ref{fig:sample-data} illustrates the input to
Algorithm~\ref{alg:sample-mean}.


\section{Algorithms}
Use a numbered algorithm float for pseudocode, with explicit inputs, outputs,
and line numbers.
\begin{algorithm}[tbp]
  \caption{Calculate the sample mean}\label{alg:sample-mean}
  \begin{algorithmic}[1]
    \Require Observations $x_1,\ldots,x_n$, where $n>0$
    \Ensure The sample mean $\bar{x}$
    \State $s\gets 0$
    \For{$i\gets 1,\ldots,n$}
      \State $s \gets s + x_i$
    \EndFor
    \State \Return $s/n$
  \end{algorithmic}
\end{algorithm}


Algorithm~\ref{alg:bootstrap-mean} illustrates a resampling method rather
than a simple arithmetic calculation. The percentile interval uses quantiles
of the bootstrap means. Its adequacy depends on the sampling setting and
should be assessed rather than assumed \citep{wasserman2004all}.
\begin{algorithm}[tbp]
  \caption{Percentile bootstrap interval for the mean}
  \label{alg:bootstrap-mean}
  \begin{algorithmic}[1]
    \Require Observations $x_1,\ldots,x_n$, $n\geq 1$, $B\geq 1$,
    and $0<\alpha<1$
    \Ensure Lower and upper percentile endpoints
    \For{$b\gets 1,\ldots,B$}
      \State Sample $n$ observations with replacement from $x_1,\ldots,x_n$
      \State Store their mean as $\bar{x}^{*(b)}$
    \EndFor
    \State \Return Empirical quantiles at $\alpha/2$ and $1-\alpha/2$
  \end{algorithmic}
\end{algorithm}


\section{Code and generated output}
Appendix~\ref{app:r-code} contains executable R chunks and their output,
including an implementation of Algorithm~\ref{alg:bootstrap-mean}. The chunks
are maintained in an Rnw file and rendered with \texttt{knitr}; the resulting
LaTeX is included in the appendix. Keep longer computational details there
and use the main chapters to explain methods and interpret results.


\section{Reference formats}
Examples include a journal article \citep{morris2019using}, a book
\citep{wasserman2004all}, a book chapter \citep{neal2011mcmc}, an edited
volume
\citep{brooks2011handbook}, a technical report \citep{breiman1999random}, a
conference
paper \citep{vaswani2017attentiona}, an arXiv preprint
\citep{vaswani2017attentionb}, a newsletter issue \citep{rnews2001news}, a
newsletter
article \citep{hornik2001editorial}, a manual \citep{venables2026introduction},
a website
\citep{rfoundationndproject}, a wiki page \citep{wikipediandbibtex}, and a
dissertation \citep{efron1964problems}. These illustrate citation formats
rather
than a research literature review. The conference and arXiv entries describe
versions of the same work; normally cite the version actually used. For
changing web pages, record the access date; for wiki pages, use a permanent
revision link in final work.
